\subsection{Actor-Critic Learner}

The learner runs the following loop inside the learned model.  An imagined
rollout is generated by the actor, the world model, and a vehicle wrapper;
each imagined step produces a per-vehicle travel-time reward and a
network-level global reward; two centralized critics estimate the
corresponding returns; and the actor is updated from their combined
advantage.  This dense credit is necessary because the global mean travel
time in Definition~7 is an episode-level objective that is available only
after a simulator episode is completed.  The per-vehicle signal is the
primary actor-learning signal, while the global critic supplies an auxiliary
aggregate value signal.
% 中文注释：学习循环：actor、world model 和 vehicle wrapper 生成 imagined rollout；每个 imagined step 产生 per-vehicle travel-time reward 和 network-level global reward；两个 centralized critics 估计对应 returns;actor 用组合 advantage 更新。需要 dense credit 是因为 Definition~7 的 global mean travel time 是 episode 级目标，只有 simulator episode 完成后才能获得。per-vehicle signal 是主要 actor-learning signal,global critic 提供辅助的 aggregate value signal。

\subsubsection{Vehicle-Level Imagination}

Imagined rollouts start from a replay state: we sample a replay episode and a
decision epoch \(t_{\mathrm{start}}\), encode the aligned traffic field as
\(Z_{t_{\mathrm{start}}}=f_\phi(X_{t_{\mathrm{start}}},G)\), and initialize
each vehicle's current segment, destination, route suffix, and arrival status
from the same epoch.  At each imagined epoch, the actor selects a route and
the decoder recovers the segment speed \(\widehat v_{t,u}\), giving the
predicted segment travel time
\begin{equation}
\begin{array}{rcl}
\widehat\tau_t(u)
&=&
\displaystyle
\frac{\ell(u)}{\widehat v_{t,u}}.
\end{array}
\end{equation}
During a decision interval \(\Delta t\), a deterministic vehicle wrapper
advances each vehicle along its selected route by consuming these travel times
in route order, while the learned transition predicts \(Z_{t+1}\).  Because
the wrapper and the transition consume the same decoded readout, vehicle
positions and the segment traffic field evolve from a single prediction rather
than two independent state estimates.  Vehicles reaching their departure time
are added at the next epoch and arrived vehicles are removed; the remaining
CAVs form \(\mathcal{I}_{t+1}^{\mathrm{im}}\) and alone contribute actor
samples, rewards, and bootstrap values.  This wrapper keeps vehicle identity
explicit without requiring the graph-level world model to predict per-vehicle
terminal variables.
% 中文注释：imagined rollout 从 replay state 初始化：采样 replay episode 和 decision epoch \(t_{\mathrm{start}}\)，编码 \(Z_{t_{\mathrm{start}}}=f_\phi(X_{t_{\mathrm{start}}},G)\)，并从同一 epoch 初始化每辆车的 current segment、destination、route suffix 和 arrival status。每个 imagined epoch,actor 选择 route,decoder 给出 segment speed \(\widehat v_{t,u}\) 和 travel time \(\widehat\tau_t(u)=\ell(u)/\widehat v_{t,u}\)。deterministic vehicle wrapper 在 \(\Delta t\) 内按 route 顺序消耗这些 travel time 推进车辆，同时 transition 预测 \(Z_{t+1}\);wrapper 与 transition 共享同一 decoded readout，因此车辆位置与 segment field 来自同一个预测。到点车辆加入、到达车辆移除，剩余 CAV 构成 \(\mathcal{I}_{t+1}^{\mathrm{im}}\)。该 wrapper 在不要求 world model 预测 per-vehicle terminal variables 的情况下显式保留 vehicle identity。

\subsubsection{Imagined Rewards}

For the route selected by CAV \(i\), let \(P_{i,t}^{\mathrm{rem}}\) denote its
suffix road segments.  The per-vehicle reward charges the imagined travel time
accumulated before the next decision or arrival:
\begin{equation}
\begin{array}{rcl}
\widehat T_{i,t}^{\mathrm{rem}}
&=&
\displaystyle
\sum_{u\in P_{i,t}^{\mathrm{rem}}}
\widehat\tau_{t+1}(u),\\[0.35em]
\delta_{i,t}
&=&
\min(\Delta,\widehat T_{i,t}^{\mathrm{rem}}),\\
r_{i,t}^{\mathrm{veh}}
&=&
-\delta_{i,t}/\kappa_r,
\end{array}
\end{equation}
where \(\kappa_r\) controls the reward scale.
% 中文注释：per-vehicle reward 由 \(\delta_{i,t}=\min(\Delta,\widehat T_{i,t}^{\mathrm{rem}})\) 和 \(r_{i,t}^{\mathrm{veh}}=-\delta_{i,t}/\kappa_r\) 给出，其中 \(\widehat T_{i,t}^{\mathrm{rem}}\) 是所选 route 剩余 suffix 的 predicted travel time。

The global reward charges each decision interval by the travel time
accumulated by all vehicles present in the network:
\begin{equation}
\begin{array}{rcl}
r_t^{\mathrm{global}}
&=&
\displaystyle
-\kappa_g\,
\frac{N_{\mathrm{active}}(t)\,\Delta t}{N_{\mathrm{veh}}},
\end{array}
\end{equation}
where \(N_{\mathrm{veh}}\) is the episode vehicle population and \(\kappa_g\)
controls the reward scale.  On real transitions, \(N_{\mathrm{active}}(t)\) is
read directly from the simulator state.  In imagined rollouts, it is computed
from the decoded count field, \(N_{\mathrm{active}}(t)=\sum_{u\in N}\widehat
n_{t,u}\) with \(\widehat n_{t,u}=\kappa(u)\hat c_{t,u}\).  Since \(\sum_t
N_{\mathrm{active}}(t)\Delta t\) is the total travel time accumulated over an
episode, the sum of \(r_t^{\mathrm{global}}\) equals \(-\kappa_g\) times the
global mean travel time of Definition~7.  Both rewards are therefore computed
from the decoded readouts rather than learned.
% 中文注释：global reward 按当前网络中所有车辆累计的通行时间对每个 decision interval 计费：\(r_t^{\mathrm{global}}=-\kappa_g N_{\mathrm{active}}(t)\Delta t/N_{\mathrm{veh}}\)。真实 transition 上 \(N_{\mathrm{active}}(t)\) 直接从 simulator 读取；imagined rollout 中用 decoded count field 求和 \(\sum_u\widehat n_{t,u}\),\(\widehat n_{t,u}=\kappa(u)\hat c_{t,u}\)。由于 \(\sum_t N_{\mathrm{active}}(t)\Delta t\) 是整个 episode 累计的总通行时间，\(r_t^{\mathrm{global}}\) 的求和等于 \(-\kappa_g\) 倍的 Definition~7 global mean travel time。两种 reward 都由 decoded readouts 计算得到，不需要学习。

\subsubsection{Critics}

The per-vehicle critic first encodes every observation independently with a
parameter-shared multilayer perceptron (MLP) \(f_\eta^{\mathrm{enc}}\), and then
applies multi-head self-attention within the active CAV set:
\begin{equation}
\begin{array}{rcl}
e_{i,t} &=& f_\eta^{\mathrm{enc}}(o_{i,t}),\\
(\widetilde e_{i,t})_{i\in\mathcal{I}_t}
&=&
\mathrm{Attention}_\eta((e_{i,t})_{i\in\mathcal{I}_t}),\\
v_{i,t}^{\mathrm{veh}}
&=&
w_v^\top\widetilde e_{i,t}+b_v.
\end{array}
\end{equation}
The attention layer lets each value estimate account for simultaneous CAV
decisions and route overlap, while \(v_{i,t}^{\mathrm{veh}}\) remains a
per-vehicle estimate of the future return.  It is trained by lambda returns
linked by vehicle identity: only observations of the same CAV are connected by
the recursion, and bootstrapping stops once the vehicle leaves the active set
(Appendix~\ref{app:math-details}).
% 中文注释：per-vehicle critic 先用 parameter-shared MLP 独立编码每个 observation，再在 active CAV set 内做 multi-head self-attention，使每个 value 能考虑 simultaneous decisions 和 route overlap，但仍为每辆车输出一个 value。它用按 vehicle identity 连接的 lambda returns 训练；车辆离开 active set 后停止 bootstrap。

The auxiliary global critic reuses the graph representation \(\bar z_t\):
\begin{equation}
\begin{array}{rcl}
v_t^{\mathrm{global}}
&=&
V_\omega^{\mathrm{global}}
(\bar z_t).
\end{array}
\end{equation}
The planned CAV pressure \(\Psi_t^{\mathrm{CAV}}\) conditions only the
world-model transition and never enters a value baseline, so
\(v_t^{\mathrm{global}}\) is an action-independent estimate of the
network-level state value.  The global critic captures aggregate congestion
effects that cannot be assigned to one CAV observation alone.  It is trained
from the corresponding epoch-level lambda return
(Appendix~\ref{app:math-details}), and its detached residual enters every
active CAV's actor advantage below.
% 中文注释：辅助 global critic 复用 graph representation \(\bar z_t\)：\(v_t^{\mathrm{global}}=V_\omega^{\mathrm{global}}(\bar z_t)\)。planned CAV pressure 只作为 world-model transition 的条件，不进入任何 value baseline，因此 \(v_t^{\mathrm{global}}\) 是 action-independent 的 network-level state value 估计。global critic 用 epoch 级 lambda return 训练，其 detached residual 进入下文每辆 active CAV 的 actor advantage。

\subsubsection{Route Actor and Policy Update}

The actor scores every candidate route from a fixed-dimensional route
representation \(\varphi_{i,t,k}\), obtained by mean pooling the segment
latents along the route, and a shared context \(g_t\) containing the pooled
graph latent, the normalized decision time, and the active-CAV fraction
(Appendix~\ref{app:math-details}).  The learned policy must select one
executable route at deployment, so the scores are normalized into a
categorical distribution over the valid candidates:
\begin{equation}
\begin{array}{rcl}
\pi_\theta(k\mid o_{i,t})
&=&
\mathrm{Softmax}_{k:m_{i,t,k}=1}
\!\left(f_\theta^{\mathrm{route}}(\varphi_{i,t,k},g_t)\right).
\end{array}
\end{equation}
Candidates with \(m_{i,t,k}=0\) are excluded from the softmax, and invalid
slots are masked out when fewer than \(K\) candidates exist.  Parameter
sharing across CAVs and candidate slots keeps the policy applicable to a
variable number of active vehicles.
% 中文注释：actor 从 mean-pooled route representation \(\varphi_{i,t,k}\) 和 shared context \(g_t\) 为每条候选打分，并对有效候选做 masked softmax:\(\pi_\theta(k\mid o_{i,t})=\mathrm{Softmax}_{k:m_{i,t,k}=1}(f_\theta^{\mathrm{route}}(\varphi_{i,t,k},g_t))\)。\(m_{i,t,k}=0\) 的候选不参与 softmax；候选不足 \(K\) 时无效 slot 被 mask。CAV 与候选 slot 之间的 parameter sharing 使 policy 适用于可变数量的 active vehicles。

For policy optimization, let \(\pi_{\mathrm{old}}\) denote the behavior policy
that produced \(a_{i,t}\).  The detached per-vehicle and global advantages and
their actor combination are
\begin{equation}
\begin{array}{rcl}
A_{i,t}^{\mathrm{veh}}
&=&
\mathrm{sg}\!\left(
G_{i,t}^{\lambda}-v_{i,t}^{\mathrm{veh}}
\right),\\
A_t^{\mathrm{global}}
&=&
\mathrm{sg}\!\left(
G_t^{\mathrm{global}}-v_t^{\mathrm{global}}
\right),\\
A_{i,t}^{\mathrm{actor}}
&=&
A_{i,t}^{\mathrm{veh}}+\alpha_g A_t^{\mathrm{global}},
\end{array}
\end{equation}
where \(G_{i,t}^{\lambda}\) and \(G_t^{\mathrm{global}}\) are the lambda-return
targets of the two critics, \(\mathrm{sg}\) stops gradients through the return
targets, and \(\alpha_g\) weights the global credit shared by CAV actions from
the same imagined epoch, so policy improvement internalizes network-wide
congestion while retaining vehicle-specific credit.  The actor is optimized
with the standard clipped PPO surrogate with entropy regularization
(Appendix~\ref{app:math-details}).  Thus the global critic affects the actor
through \(A_t^{\mathrm{global}}\), rather than only through a critic loss.
% 中文注释：vehicle advantage 与 global advantage 分别为 \(A_{i,t}^{\mathrm{veh}}=\mathrm{sg}(G_{i,t}^{\lambda}-v_{i,t}^{\mathrm{veh}})\) 和 \(A_t^{\mathrm{global}}=\mathrm{sg}(G_t^{\mathrm{global}}-v_t^{\mathrm{global}})\),actor 使用 \(A_{i,t}^{\mathrm{actor}}=A_{i,t}^{\mathrm{veh}}+\alpha_gA_t^{\mathrm{global}}\)。\(\mathrm{sg}\) 阻断 return targets 上的梯度，\(\alpha_g\) 控制 global credit 的权重；同一 imagined epoch 的 active CAVs 共享 global term，因此 policy update 同时保留 vehicle-specific credit 并纳入 network-wide congestion credit。actor 用标准 clipped PPO surrogate 加 entropy regularization 优化(附录)。global critic 不只产生 critic loss，它的 detached lambda advantage 直接进入 actor objective。
